\section{A weaker condition based on the gap between training and test margins}
\label{sec:margin-band-variant}

The primary condition uses the score clipped at the training margins.
The argument admits a further weakening which ignores changes entirely
inside the region already counted as a test-margin error. We present
this as a mathematical variant; it does not retroactively change the
primary condition used in the empirical protocol.

\begin{theorem}[Monotone score transforms]
\label{thm:monotone-score-transform}
Fix, independently of the training sample, an odd, nondecreasing and
$1$-Lipschitz map $F:\mathbb R\to\mathbb R$ that is constant on
$[1/2,\infty)$. Put
\[
 A=F(1/2),\qquad B=F(1/4),\qquad \Delta=A-B>0.
\]
Replace the main support condition by
\begin{equation}\label{eq:transformed-score-support}
 |F(z_w(x;\eta))-F(z_w(x))|\le\Delta/4
 \quad\text{for every }x\text{ and every }
       \|\eta\|_*\le r_c(w).
\end{equation}
Keep the same Gaussian prior, observed training margin $yz>1/2$,
and test error event $yz\le1/4$. Then the main all-width, all-$\gamma\ge1$
polynomial sample guarantee remains valid, with additional polynomial
dependence on $\Delta^{-1}$.
The same replacement is valid in the controlled-exception localized
prior theorem, with its $\rho+\varepsilon$ conclusion.
\end{theorem}

\begin{proof}
The existing target cylinder has clean and perturbed signed output
larger than one, so all these outputs have transformed score $fA$.
Its support membership and raw Gaussian probability are unchanged.
The readout-tail and entropy estimates are also unchanged.
After the readout truncation $E_v\le B_v$, use common compression radius
\[
 r=\min\{1,2\Delta/3,c/\sqrt{1+HB_v}\}.
\]
On a successful coupling, $1$-Lipschitzness of $F$ gives transformed
score discrepancy at most
\[
 \Delta/4+r/8\le\Delta/3.
\]
At a training input, $yF(z_w)=A$, so the compressed candidate has
signed transformed score at least $A-\Delta/3$. At a test-margin
error, monotonicity and oddness give
$yF(z_w)=F(yz_w)\le B$, so its compressed signed score is at most
$B+\Delta/3$. These endpoints are separated by $\Delta/3$.
Use the $[0,1]$ ramp equal to one below $B+\Delta/3$, zero above
$A-\Delta/3$, and linear between them. Its slope magnitude is
$3/\Delta$. The same finite-class bounded-loss concentration and
coupling argument therefore gives the same risk inequality.
All common class parameters are polynomial in the old parameters and
$\Delta^{-1}$. Masking the ramp outside a controlled exceptional set
also gives exactly the earlier localized proof.
\end{proof}

\paragraph{The margin-band score.}
The particularly natural choice is
\begin{equation}\label{eq:margin-band-score}
 U(z)=\operatorname{sign}(z)
          \min\{1/4,(|z|-1/4)_+\}.
\end{equation}
It is zero on $[-1/4,1/4]$, grows linearly toward the training margins,
and saturates at $\pm1/4$ once $z$ reaches $\pm1/2$.
Here $A=1/4$, $B=0$ and $\Delta=1/4$. The support tolerance is
still $1/16$, and the common compression radius still satisfies
$r\le1/6$. The separating ramp thresholds are $1/12$ and $1/6$.
Thus all constants in the main statistical argument remain unchanged.

The score measures signed excess beyond the test-error threshold.
For example, changing an output from $0.1$ to $-0.1$ leaves both
transformed scores zero. Both outputs already fail the $1/4$ margin
criterion for either label; forbidding this
fluctuation is unnecessary for the proof. In contrast, movement from
an unconfident output into the training-confidence region remains
controlled. Training-margin satisfaction and posterior probability of
correct test answers retain their original meanings; the transform
only defines the prior's stability condition.

This is genuinely weaker than clipping at $\pm1/2$: with
$T(z)=\operatorname{clip}(z,-1/2,1/2)$,
\[
 U(z)=\operatorname{sign}(T(z))(|T(z)|-1/4)_+,
\]
and the map on the right is $1$-Lipschitz in $T(z)$.
Every network satisfying the primary condition therefore satisfies the
margin-band condition at the same radius and tolerance.

There is a precise minimality statement within this class of score
transforms. Suppose $F$ satisfies the theorem and retains the full
available gap $F(1/2)-F(1/4)=1/4$. Monotonicity and
$1$-Lipschitzness imply, with $B=F(1/4)\ge0$,
\[
 F(z)=B+z-1/4\quad(1/4\le z\le1/2),\qquad
 0\le F(z)\le B\quad(0\le z\le1/4).
\]
Using saturation and oddness elsewhere gives
\[
 U(z)=\operatorname{sign}(F(z))(|F(z)|-B)_+.
\]
Thus stability of any such $F$ implies stability of $U$ at the same
radius and tolerance. In this explicitly specified class, the
margin-band score imposes the weakest stability requirement while
retaining the full training-to-test gap. This is not a claim of
optimality among all possible prior conditions.

A raw-output interval certificate $z_w(x;\eta)\in[z-B_x,z+B_x]$
certifies this variant whenever
\[
 \max\{U(z+B_x)-U(z),U(z)-U(z-B_x)\}\le1/16.
\]
Existing interval-certified attacks can likewise be evaluated with
$U$ without a new attack search. A failure or success of the primary
clipped-score condition need not have the same outcome for this variant.
